% from ega4-2-fr.tex
\oldpage[IV-2]{218}

\item[\emph{(ii)}] Si $X$ est réduit, son normalisé $X'$ (\textbf{II}, 6.3.8) est fini
sur $X$.
\item[\emph{(iii)}] Les ensembles $\operatorname{Reg}(X)$, $\operatorname{Nor}(X)$,
$\mathrm U_{\mathrm R_n}(X)$ sont ouverts dans $X$, ainsi que
$\mathrm U_{\mathrm C_n}(\mathcal F)$, $\mathrm U_{\mathrm S_n}(\mathcal F)$
et $\operatorname{CM}(\mathcal F)$ pour tout $\mathcal O_X$-Module cohérent
$\mathcal F$.
\end{enumerate}
\end{proposition}

\begin{proof}
Cela résulte aussitôt de (7.8.3, (ii), (vi) et (iv)). Notons aussi que (7.8.3,
(ix)) est valable sans modification d'énoncé lorsque $X$ est un préschéma
excellent.
\end{proof}

\subsection{Anneaux excellents et résolution des singularités}
\label{subsection:IV.2.7.9-fr}

\begin{env}[7.9.1]
\label{IV.2.s7.9.1-fr}
Étant donné un préschéma \emph{localement noethérien réduit} $X$, on appelle
\emph{morphisme résolvant} pour $X$ un morphisme $f:X'\to X$ tel que $X'$ soit
\emph{régulier} et que $f$ soit \emph{propre} et \emph{birationnel}. Lorsqu'il
existe un tel morphisme, on dit qu'on peut \emph{résoudre les singularités} de
$X$ (ou plus simplement que l'on peut \emph{résoudre} $X$). Par exemple, si
$A$ est un \emph{anneau japonais} de dimension $1$, on peut résoudre
$X=\operatorname{Spec}(A)$ puisque le morphisme $X'\to X$ (où $X'$ est le
\emph{normalisé} de $X$) est fini (donc propre) et birationnel, et que les
anneaux locaux de $X'$ sont des anneaux intégralement clos de dimension $1$,
donc des anneaux de valuation discrète, et par suite sont réguliers.
\end{env}

\begin{env}[7.9.2]
\label{IV.2.s7.9.2-fr}
Il est clair que si l'on peut résoudre $X$, on peut aussi résoudre tout
préschéma induit sur un \emph{ouvert} de $X$ et tout préschéma local
$\operatorname{Spec}(\mathcal O_x)$ de $X$ (\textbf{II}, 5.4.2). D'autre part,
si l'on peut résoudre $X$, il est clair qu'il existe dans $X$ un ouvert
\emph{partout} dense contenu dans $\operatorname{Reg}(X)$. Ces remarques
montrent aussitôt que si, pour tout sous-préschéma \emph{fermé intègre} $Y$ de
$X$ et tout $Y$-préschéma $Y'$ \emph{intègre, fini} et \emph{radiciel} sur
$Y$, on peut résoudre $Y'$, alors les ouverts affines de $X$ \emph{vérifient
la condition} (iii) de (7.8.2).

D'autre part, pour les schémas locaux, on a la proposition suivante :
\end{env}

\begin{proposition}[7.9.3]
\label{IV.2.s7.9.3-fr}
Soit $A$ un anneau local noethérien réduit, et supposons que l'on puisse
résoudre $\operatorname{Spec}(A)$. Alors les fibres du morphisme canonique
$\operatorname{Spec}(\widehat A)\to\operatorname{Spec}(A)$ aux points maximaux
de $\operatorname{Spec}(A)$ sont régulières.
\end{proposition}

\begin{proof}
Posons $X=\operatorname{Spec}(A)$, $X'=\operatorname{Spec}(\widehat A)$, et
soit $g:X'\to X$ le morphisme canonique. Soit $f:Y\to X$ un morphisme
résolvant; posons $Y'=X'\times_XY$, et soient $f':Y'\to X'$ et $g':Y'\to Y$
les projections canoniques, de sorte que l'on a un diagramme commutatif
\[
\xymatrix{
X'\ar[d]_g & Y'\ar[l]_{f'}\ar[d]^{g'}\\
X & Y\ar[l]^f
}
\]

Il suffira de prouver que le préschéma $Y'$ est \emph{régulier} : en effet,
comme $f$ est birationnel et de type fini, il y a un ouvert $U$ partout dense
dans $X$ tel que la restriction $f^{-1}(U)\to U$ de $f$ soit un isomorphisme
et que $f^{-1}(U)$ soit partout dense dans $Y$ (\textbf{I}, 6.5.5); les
préschémas induits respectivement sur l'ouvert $g^{-1}(U)$ de $X'$ et l'ouvert
$g'^{-1}(f^{-1}(U))$ de $Y'$ sont donc isomorphes; cela prouve que
$g^{-1}(U)$ est régulier, et \emph{a fortiori} il en est de même des fibres de
$g$ aux points maximaux de $X$, qui sont contenus dans $U$.

